% TeX ignores anything on a line after a percent sign %
% The next two lines define fonts for the title
\font\xmplbx = cmbx10 scaled \magstep1
\font\xmplbxti = cmbxti10 scaled \magstep1
% Now here's the title.
\leftline{\xmplbx Example 1:\quad\xmplbxti Entering simple text}
\vglue 0.5\baselineskip % skip an extra half line

The definitive source of information on \TeX\ is {\sl The \TeX book}, by D.~E.~Knuth, author of the
\TeX\ program. {\sl\TeX\ for the Impatient\/} by P.~W.~Abrahams, {\it et al.}\ is also an excellent
reference.  The \TeX\ software is free; you can obtain an MS-DOS version, e.g., on the NUS ftp site
({\tt ftp.nus.sg}) in the subdirectory {\tt /pub/zf/TeX/systems/msdos/emtex}.

It's easy to prepare ordinary text for \TeX\ since \TeX\ usually doesn't care about how you break
up your input into lines.  It treats the end of a line of text like a space.% 
\footnote \dag{\TeX\ treats a tab like a space too, as we point out in this {\it footnote}.} 
If you don't want a space there, put a per% 
cent sign (the comment character) at the end of the line.  \TeX\ ignores spaces
at the start of a line, and treats more than one space as equivalent to a single space, even after
a period.  You indicate a new paragraph by skipping a line (or more than one line).

When \TeX\ sees a period followed by a space (or the end of the line, which is equivalent), it
ordinarily assumes you've ended a sentence and inserts a little extra space after the period.  It
treats question marks and exclamation points the same way.

But \TeX's rules for handling periods sometimes need fine tuning.  \TeX\ assumes that a capital
letter before a period doesn't end the sentence, so you have to do something a little different if,
say, you're writing about DNA\null.  It's a good idea to tie words together in references such as
``see Fig.~8'' and in names such as V.~I. Lenin and in $\ldots$ so that \TeX\ will never split them
in an awkward place between two lines.  (The three dots indicate an ellipsis.)

You should put quotations in pairs of left and right single ``quotes'' so that you get the correct
left and right double quotations marks.  ``For adjacent single and double quotation marks, insert
a `thinspace'\thinspace''.  You can get en-dashes--like so, and em-dashes---like so.
